% ============================================================================
%  JobScope — LaTeX project documentation
%  Compile with:  pdflatex JobScope_Guide.tex   (run twice for the ToC)
%  Packages required (all standard): tikz, listings, booktabs, hyperref,
%  enumitem, geometry, xcolor, titlesec.
% ============================================================================
\documentclass[11pt,a4paper]{article}

% ---- packages --------------------------------------------------------------
\usepackage[utf8]{inputenc}
\usepackage[T1]{fontenc}
\usepackage{amssymb}
\usepackage[margin=2.6cm]{geometry}
\usepackage{xcolor}
\usepackage{listings}
\usepackage{booktabs}
\usepackage{enumitem}
\usepackage{tikz}
\usetikzlibrary{arrows.meta,positioning,shadows.blur}
\usepackage{titlesec}

% ---- colours (defined BEFORE hyperref so its colour options resolve) -------
\definecolor{indigo}{HTML}{4F46E5}
\definecolor{violet}{HTML}{7C3AED}
\definecolor{surface}{HTML}{F1F5F9}
\definecolor{codebg}{HTML}{0F172A}
\definecolor{teal}{HTML}{0D9488}

\usepackage[colorlinks=true,linkcolor=indigo!70,urlcolor=teal]{hyperref}

% ---- code listing style ------------------------------------------------------
\lstset{
  basicstyle=\ttfamily\footnotesize\color{white},
  backgroundcolor=\color{codebg},
  commentstyle=\color{teal},
  keywordstyle=\color{orange!90},
  stringstyle=\color{green!70},
  numbers=left, numberstyle=\tiny\color{gray}, numbersep=6pt,
  frame=single, rulecolor=\color{gray!40}, breaklines=true,
  showstringspaces=false, aboveskip=10pt, belowskip=10pt,
}

% ---- section styling ---------------------------------------------------------
\titleformat{\section}{\Large\bfseries\color{indigo}}{\thesection}{0.6em}{}
\titleformat{\subsection}{\large\bfseries\color{violet}}{\thesubsection}{0.5em}{}
\setlist{itemsep=2pt, topsep=4pt}

% ---- handy commands ------------------------------------------------------------
\newcommand{\file}[1]{\texttt{\small #1}}
\newcommand{\tool}[1]{\texttt{#1}}
\newcommand{\nodepill}[1]{{\footnotesize\bfseries\color{white}\colorbox{indigo}{#1}}}
\newcommand{\link}[1]{\url{#1}}

\setlength{\parindent}{0pt}
\setlength{\parskip}{6pt}

\begin{document}

% ============================================================================
%  TITLE PAGE
% ============================================================================
\begin{titlepage}
  \centering
  \vspace*{2.5cm}
  {\color{indigo}\rule{\textwidth}{2pt}}

  \vspace{0.9cm}
  {\fontsize{34}{40}\selectfont\bfseries\color{indigo} JobScope}\\[6pt]
  {\Large A LangGraph\,$\cdot$\,Ollama\,$\cdot$\,MCP job-hunting agent with a live dashboard}

  \vspace{0.9cm}
  {\color{indigo}\rule{\textwidth}{2pt}}

  \vspace{2.6cm}
  \begin{minipage}{0.78\textwidth}
    \small
    \textbf{Abstract.} JobScope is a locally-run job-hunting agent.
    A LangGraph actor turns a target role and the user's resume skills into
    keyword queries, dispatches them to several job boards through MCP tool
    servers, and \emph{deterministically} falls back to Gmail job-alert
    emails when the boards are bot-blocked. A FastAPI dashboard streams the
    whole run to the browser as Server-Sent Events, renders listings as an
    Apply-able flashcard deck, and animates a circuit-style pipeline with
    per-step hover logs. Everything runs offline against a local LLM (Ollama);
    no external API keys beyond a free GitHub PAT and a Gmail app password.
  \end{minipage}

  \vspace{3cm}
  {\small \today \\[4pt]
   Build target: Python 3.14, Windows / PowerShell \\[2pt]
   File: \file{docs/JobScope\_Guide.tex}}
\end{titlepage}

\tableofcontents
\clearpage

% ============================================================================
\section{Introduction}
\label{sec:intro}

JobScope is a job-hunting assistant that the user runs \emph{entirely on their
own machine}. It is the final result of a seven-step tutorial that grows, step
by step, from a bare LangGraph state machine (\file{step1}) to a full agent
with multiple MCP tool servers and a deterministic fallback
(\file{agent.py}).

\begin{itemize}
  \item \textbf{Agent framework:} LangGraph (StateGraph, prebuilt \file{ToolNode},
        \file{tools\_condition} router).
  \item \textbf{Model:} Ollama running locally on \link{http://localhost:11434}
        (default \tool{llama3.1}, configurable e.g.\ \tool{qwen2.5}). No paid
        API.
  \item \textbf{Tool protocol:} Model Context Protocol (MCP) over stdio — each
        server is a Python subprocess that the agent's client launches.
  \item \textbf{Job sources:} nine job boards via a multi-board scraper, plus a
        Gmail IMAP fallback.
  \item \textbf{Constraint:} the user lives in India and only wants
        \textbf{remote} jobs that are \textbf{not} based in India. This is
        enforced as deterministic code, not prose.
\end{itemize}

% ============================================================================
\section{Architecture overview}
\label{sec:arch}

\newcommand{\archnode}[4]{%
  \node[draw=#2!70, fill=#2!12, rounded corners=4pt,
        minimum width=#4, minimum height=2.1em, font=\small,
        blur shadow={shadow blur steps=5}] (#1) {#3};
}

\begin{center}
\begin{tikzpicture}[>=Stealth, node distance=1.15cm and 1.6cm]
  \archnode{dash}{indigo}{FastAPI dashboard\\\scriptsize (dashboard.py)}{3.4cm}
  \archnode{agent}{violet}{LangGraph agent\\\scriptsize (agent.py)}{3.0cm}
  \archnode{mcp}{teal}{MCP servers (stdio subprocesses)\\ \scriptsize github \textbar\ scraper \textbar\ gmail}{4.4cm}
  \archnode{boards}{gray!70}{Job boards (Ind/Lkd/Nk/GD/Fit/IS/WWR/Remo/Arb)\\ \scriptsize WWR+Remo+Arb+IS-WFH live $\Rightarrow$ mixed deck; big 6 walled $\Rightarrow$ fallback}{4.6cm}
  \archnode{mail}{gray!70}{Gmail IMAP job alerts}{3.2cm}
  \archnode{ollama}{violet!60}{Ollama LLM\\ \scriptsize localhost:11434}{2.8cm}

  \draw[->, thick, indigo] (dash) -- node[above,font=\scriptsize]{SSE on port 8000} (agent);
  \draw[->, thick, violet] (agent) -- node[above,font=\scriptsize]{tool calls} (mcp);
  \draw[->, thick, gray]   (mcp)  -- node[below,font=\scriptsize]{HTTPS} (boards);
  \draw[->, thick, gray]   (mcp)  -- node[below,font=\scriptsize]{IMAPS:993} (mail);
  \draw[->, thick, violet!60] (agent) -- node[right,font=\scriptsize]{prompt} (ollama);
  \draw[->, thick, teal]   (boards) -- node[right,font=\scriptsize]{listings or 403/CAPTCHA} (mcp);
  \draw[->, thick, teal]   (mail)   -- node[right,font=\scriptsize]{alert emails} (mcp);
  \draw[->, dashed, thick, violet] (mcp) to[bend right=25] node[below,font=\scriptsize]{fallback route} (mail);
\end{tikzpicture}
\end{center}

\begin{itemize}
  \item The dashboard and the chat CLI share the \emph{same} compiled graph —
        \file{create\_agent()} is the single source of truth.
  \item MCP servers are subprocesses with \tool{stdio} transport; they are
        \emph{not} network services (only ports 8000 and 11434 are open locally).
  \item When \emph{all} boards are blocked, a deterministic graph edge routes
        the run to a fallback node that queries Gmail instead.
\end{itemize}

% ============================================================================
\section{Components and project layout}
\label{sec:files}

\begin{center}
\begin{tabular}{@{}llp{7.4cm}@{}}
\toprule
\textbf{File} & \textbf{Role} & \textbf{Purpose}\\
\midrule
\file{agent.py} & core & shared agent: config, MCP client, graph, router, filters\\
\file{dashboard.py} & backend & FastAPI server, SSE streaming, resume API, JSONL logging\\
\file{static/index.html} & frontend & job dashboard: pipeline, flashcards, drawers\\
\file{mcp\_server\_github.py} & MCP & GitHub repo / README tools\\
\file{mcp\_server\_indeed\_scraper.py} & MCP & multi-board \tool{search\_job\_boards} scraper\\
\file{mcp\_server\_gmail.py} & MCP & \tool{search\_job\_emails} via IMAP + app password\\
\file{resume\_generator.py} & core & per-job LaTeX resume builder + HTML preview\\
\file{resume\_data.py} & data & canonical resume content (skills/experience/education)\\
\file{ui.py} & frontend & Chainlit chat UI (alternative frontend)\\
\file{step7\_scraper\_with\_fallback.py} & CLI & command-line full run with fallback\\
\file{step1..step6\_*.py} & tutorial & seven-step learning progression\\
\file{requirements.txt} & setup & pinned Python dependencies\\
\file{.env.example} & setup & secret/config template\\
\file{logs/} & output & \file{runs-YYYY-MM-DD.jsonl} + \file{errors.jsonl}\\
\bottomrule
\end{tabular}
\end{center}

\subsection{The seven-step tutorial progression}
The \file{step*.py} files build the concept understanding:
\begin{enumerate}
  \item \file{step1\_basics.py} — bare LangGraph state machine with three nodes.
  \item \file{step2\_agent\_with\_tool.py} — LLM plus a plain Python tool.
  \item \file{step3\_github\_tool.py} — decorated \tool{@tool} functions hitting GitHub.
  \item \file{step4\_mcp\_client.py} — client-side MCP tool discovery.
  \item \file{step5\_two\_mcp\_servers.py} — two servers at once.
  \item \file{step6\_gmail\_source.py} — Gmail as a job source.
  \item \file{step7\_scraper\_with\_fallback.py} — full agent wiring the scraper
        and the Gmail fallback.
\end{enumerate}

% ============================================================================
\section{The LangGraph agent (\file{agent.py})}
\label{sec:agent}

\subsection{State}
The graph uses \tool{AgentState}(\file{MessagesState}) with one extra key:
\tool{fallback\_used} — a flag that stops the router from looping
\emph{scrape-blocked \(\to\) fallback \(\to\) scrape-blocked}.

\subsection{Nodes and edges}
\begin{itemize}
  \item \nodepill{agent} --- on a fresh hunt this \emph{skips the LLM}: it emits
        one \tool{AIMessage} whose parallel tool\_calls fan out to
        \tool{search\_job\_boards} and to both Gmail alert senders
        simultaneously (zero inference). Afterwards the final answer is built
        \textbf{in code by default} (\tool{AGENT\_SUMMARY\_MODE=template}) ---
        instant, every title/company/location/Apply link preserved from the
        structured payload. With \tool{AGENT\_SUMMARY\_MODE=llm} the Ollama LLM
        writes it instead --- an \emph{unbound} instance on a trimmed context
        (raw tool dumps are removed from state first), to keep the expensive
        CPU inference short.\par\nopagebreak
  \item \nodepill{tools} --- prebuilt \tool{ToolNode} executes the MCP tools;
        the parallel tool\_calls run concurrently.
  \item \nodepill{prep} --- picks the winning source in code: boards if they
        returned any listings, otherwise the filtered Gmail fallback; then
        hands the summary LLM one self-contained payload (with the
        \tool{###JOBS\_JSON###} feed for flashcards).
  \item \nodepill{fallback} --- legacy sequential path (LLM-decided runs where
        the scraper was blocked); pulls both alert senders from Gmail in
        parallel.
\end{itemize}

Edges:
\begin{lstlisting}[language=Python]
START -> agent -> (tools_condition)
  tools_condition: "tools"  -> tools
                   "__end__"-> END
tools -> route_after_tools: "prep"    (parallel collect ran)
                            "fallback"(scraper was blocked, legacy)
                            "agent"   (summarise normally)
prep -> agent
fallback -> agent
\end{lstlisting}

The router \tool{route\_after\_tools} routes to \tool{prep} whenever Gmail
tool results are already in state (the parallel-collect path), else to the
legacy \tool{fallback} when the job-board tool's text contains a blocked
marker, else back to \tool{agent} to write the final answer. \tool{fallback\_used}
stops any of this from looping.

\subsection{Blocked markers}
A multi-board run can fail on \emph{some} boards yet still return jobs, so the
literal word \tool{"blocked"} alone is \textbf{not} a fallback trigger. The
markers that are (\file{agent.py:53}):
\begin{lstlisting}
"all boards were blocked", "captcha", "http 403", "fall back",
"no job links", "couldn't extract", "request failed", "js-rendered",
\end{lstlisting}

% ============================================================================
\section{Job filtering rules}
\label{sec:filters}

All filtering is \textbf{deterministic code}, applied to fallback data and
scraper listings before the LLM sees them:

\begin{description}[leftmargin=5.5em]
  \item[Seniority] \tool{SENIORITY\_FILTER} drops titles matching
        \tool{senior\,|\,lead\,|\,manager\,|\,principal\,|\,staff\,|\,architect\,|\,director\,|\,head of}.
  \item[Experience] \tool{EXPERIENCE\_FILTER} drops lines matching
        \tool{\textbackslash b(?:[5-9]\textbar\textbackslash d\{2,\})+\textbackslash+ years/yrs}.
  \item[Geography] \tool{INDIA\_TOKENS} (\file{agent.py:73}) matches Indian
        cities/states/``Remote-India'' and, when \tool{remote\_only} is on,
        the listing is excluded.
  \item[Links] \tool{\_keep\_job} additionally requires a valid \tool{http(s)}
        link so every card is Apply-able.
\end{description}

The system prompt mirrors these rules for the LLM (so it re-rationalises them
in prose), but the graph does not rely on it.

% ============================================================================
\section{MCP tool servers}
\label{sec:mcp}

Every server declares tools with \tool{@mcp.tool()} and runs with
\tool{mcp.run(transport="stdio")}. The client
(\tool{MultiServerMCPClient}) spawns each server with \tool{sys.executable}.

\begin{center}
\begin{tabular}{@{}lll@{}}
\toprule
\textbf{Server} & \textbf{Tools} & \textbf{Notes}\\
\midrule
\file{mcp\_server\_github.py} & \tool{list\_github\_repos}, \tool{get\_repo\_readme} & GitHub API + PAT\\
\file{mcp\_server\_indeed\_scraper.py} & \tool{search\_job\_boards} & 9 boards, per-source blocked status\\
\file{mcp\_server\_gmail.py} & \tool{search\_job\_emails} & IMAP, 15\,s connect timeout\\
\bottomrule
\end{tabular}
\end{center}

\subsection{Multi-board scraper}
\tool{search\_job\_boards(query, location="Remote", remote\_only=True)} tries
Indeed, LinkedIn, Naukri, Glassdoor, Foundit, Internshala, WeWorkRemotely,
Remotive and Arbeitnow concurrently (one pool, one DNS was the hang), dedupes
results, drops senior/5+-year listings, applies the India filter (work-from-
home roles are exempted), and reports each board as \emph{a number of
listings} or \emph{blocked (reason)}. A hard deadline caps the whole hunt so
a stuck fetch can never hang the run. Most big Indian boards still bot-wall
a plain HTTP client (403/CAPTCHA or connection refusal), but WeWorkRemotely
serves static HTML and Remotive + Arbeitnow expose keyless JSON feeds, and
Internshala's work-from-home internships page scrapes cleanly — so a mixed,
multi-source deck is the norm rather than the Gmail fallback.
Output ends with a machine-readable block:

\begin{lstlisting}[language=Python]
###JOBS_JSON###
{"sources": {"Indeed": 12, ...}, "blocked": [...], "total": 12,
 "remote_only": true, "jobs": [{"source","title","company",
  "location","salary","link"}, ...]}
\end{lstlisting}

\subsection{Gmail server}
\tool{search\_job\_emails(sender, days\_back, max\_results)}:
\begin{enumerate}
  \item connects over IMAPS (15\,s timeout so a blocked network fails fast);
  \item searches \tool{FROM sender SINCE date}, newest first;
  \item strips the huge tracking-URL footers that waste LLM context;
  \item caps each body at 1200 characters;
  \item parses each email into a structured listing
        (\tool{\_extract\_listing}) keeping the one real ``View job'' URL,
        and appends its own \tool{###JOBS\_JSON###} block.
\end{enumerate}

% ============================================================================
\section{Resume pipeline and role suggestions}
\label{sec:resume}

The dashboard turns a CV (pdf/txt/md) into the agent's search terms:

\begin{enumerate}
  \item \tool{\_extract\_text} — text from PDF (\tool{pypdf}) or plain text decode.
  \item \tool{extract\_skills} — in the default \file{strict} mode it takes
        \textbf{only} keywords that literally appear in the CV: every entry of
        the ``Technical Skills'' block if one exists (colons, pipe-separated
        table rows or bullets), otherwise known skills found verbatim in the
        text. No invented/norexternal terms; \tool{RESUME\_EXTRACT\_MODE}
        \file{auto} re-enables lexicon enrichment and \file{llm} forces the
        local model.
  \item \tool{suggest\_roles} — a rubric scores role titles against the skills
        (e.g.\ \tool{sql python power bi tableau} \(\to\) \emph{Data Analyst}).
  \item \tool{infer\_location} — regex/city lookup on the CV text; the inferred
        location also drives \tool{remote\_only} (the sample CV resolved to
        ``Hyderabad, India'').
\end{enumerate}

The top role suggestion becomes the \tool{target\_role}, the full list is stored
in \tool{role\_suggestions}, and \tool{skills} feed the keyword-driven search
query (\tool{search\_query\_from} uses the three strongest skill keywords
before falling back to the plain role title).

\subsection{Per-job LaTeX resume generation}
Each flashcard's \emph{LaTeX resume} button calls \file{POST /api/resume/gen}
(\tool{resume\_generator.py}), which produces an ATS-friendly 10\,pt article
document from the user's canonical resume (\tool{resume\_data.py}):

\begin{enumerate}
  \item \tool{desc\_snippet} --- best-effort fetch of the posting page for a
        description snippet (boards are usually bot-walled; failures are
        silent and the resume is still tailored from role + skills).
  \item \tool{\_rank\_skills} --- counts how often each skill group's keywords
        appear in the posting and reorders \tool{resume\_data.SKILLS} so the
        most job-relevant groups come first; matched keywords also seed the
        Objective.
  \item \tool{build\_resume} --- emits \file{.tex} from a raw \tool{ATS\_TEMPLATE}
        (standard \verb|\section*| headings, \verb|\heading[4]| and
        \verb|\subheading[2]| macros, accent \verb|1F4E79|, no tables, no
        columns). Headers/links come from \tool{resume\_data} optionally
        overridden by the uploaded CV (\tool{extract\_contact}).
  \item \tool{tex\_to\_html} --- the same source is converted into a standalone
        HTML preview so the browser shows the resume without a TeX install.
\end{enumerate}

The dashboard returns \tool{\{ok, filename, tex, preview, desc\_fetched\}}; the
browser renders the preview, and offers \emph{Copy LaTeX} and \emph{Download
.tex} (the file compiles with a plain \tool{pdflatex}).

% ============================================================================
\section{Dashboard backend (\file{dashboard.py})}
\label{sec:dash}

\subsection{Endpoints}
\begin{center}
\begin{tabular}{@{}ll@{}}
\toprule
\textbf{Method/path} & \textbf{Purpose}\\
\midrule
\file{GET /} & serve the single-page UI\\
\file{GET /api/config} & current config + discovered tool names\\
\file{POST /api/config} & rebuild the agent with new settings\\
\file{POST /api/resume} & upload CV $\Rightarrow$ skills, suggestions, location\\
\file{POST /api/resume/gen} & build a tailored LaTeX resume + preview for one listing\\
\file{POST /api/run} & run the agent; \textbf{SSE} stream of events\\
\file{GET /api/logs} & recent run history + errors\\
\bottomrule
\end{tabular}
\end{center}

\subsection{SSE event frames}
\file{/api/run} responds \tool{text/event-stream} and emits framed JSON events:

\begin{center}
\begin{tabular}{@{}ll@{}}
\toprule
\textbf{Event type} & \textbf{Meaning}\\
\midrule
\file{status} & run life-cycle (starting / done / error)\\
\file{token} & streamed LLM token chunks\\
\file{tool\_start} / \file{tool\_end} & an MCP tool began / finished (with preview)\\
\file{jobs} & parsed \tool{###JOBS\_JSON###} payload $\Rightarrow$ flashcard deck\\
\file{fallback} & the deterministic Gmail fallback fired\\
\file{error} & fatal failure (banner with Retry/Restart)\\
\bottomrule
\end{tabular}
\end{center}

\subsection{Logging}
Every run and error is appended as JSONL:
\begin{itemize}
  \item \file{logs/runs-YYYY-MM-DD.jsonl} — run id, message, skills, tools used,
        duration, job counts, model.
  \item \file{logs/errors.jsonl} — failures with full tracebacks.
\end{itemize}

% ============================================================================
\section{Frontend (\file{static/index.html})}
\label{sec:ui}

A single HTML file (Inter + JetBrains Mono fonts with CDN fallback) that:
\begin{itemize}
  \item renders a \textbf{circuit pipeline} (Agent \(\to\) Job boards \(\to\)
        Router \(\to\) Gmail \(\to\) Report) whose connectors light up with an
        animated dash as the run progresses, and whose \textbf{hover popovers}
        show each step's live log;
  \item renders results as a \textbf{flashcard deck} — page-turn/slide
        animation, arrows, dots, swipe, keyboard — where every card carries an
        \textbf{Apply now} button opening the job URL in a new tab;
  \item offers a Settings drawer (role chips, editable skill chips, resume
        upload, remote-only toggle) and a History drawer (Retry / Restart);
  \item adds 3-D hover interactions (tilted node pills, lifted glowing
        buttons with gradient sheen, floating gradient blobs) that respect
        \texttt{prefers-reduced-motion}.
\end{itemize}

% ============================================================================
\section{Configuration}
\label{sec:config}

\subsection{\file{.env} variables}
\begin{center}
\begin{tabular}{@{}lll@{}}
\toprule
\textbf{Variable} & \textbf{Meaning} & \textbf{Default}\\
\midrule
\file{OLLAMA\_MODEL} & Ollama model name & \file{llama3.1}\\
\file{OLLAMA\_NUM\_CTX} & model context window (tokens) & \file{8192}\\
\file{TARGET\_ROLE} & role the agent matches listings to & \emph{Junior Data Analyst}\,…\\
\file{RESUME\_EXTRACT\_MODE} & \file{strict} (keywords only from the CV, no invented extras), \file{auto} (also enriches with a known-skill lexicon), \file{fast} (lexicon only), \file{llm} & \file{strict}\\
\file{AGENT\_SUMMARY\_MODE} & final answer: \file{template} (instant, in-code) or \file{llm} (local model) & \file{template}\\
\file{JOB\_DAYS\_BACK} & email search window & \file{60}\\
\file{JOB\_MAX\_RESULTS} & emails per sender & \file{10}\\
\file{GITHUB\_TOKEN} & GitHub PAT (public reads) & ---\\
\file{GITHUB\_USERNAME} & GitHub account & ---\\
\file{GMAIL\_ADDRESS} & Gmail address (IMAP) & ---\\
\file{GMAIL\_APP\_PASSWORD} & Gmail app password & ---\\
\bottomrule
\end{tabular}
\end{center}

\verb|.env| never gets committed — \file{.env.example} is the template.

\subsection{AgentConfig}
\tool{AgentConfig} (dataclass in \file{agent.py}) centralises model, context
window, role, search window/max results, \tool{skills}, \tool{resume\_text},
\tool{remote\_only}, \tool{exclude\_locations} and \tool{role\_suggestions};
\tool{load\_config()} fills it from \file{.env} with sane defaults.

% ============================================================================
\section{Setup and usage}
\label{sec:usage}

\begin{lstlisting}[language=PowerShell, numbers=none]
python -m venv .venv
.\.venv\Scripts\Activate.ps1
.\.venv\Scripts\python.exe -m pip install -r requirements.txt

Copy-Item .env.example .env   # fill in YOUR tokens
ollama pull llama3.1          # https://ollama.com/download
\end{lstlisting}

Launch options:
\begin{center}
\begin{tabular}{@{}ll@{}}
\toprule
\textbf{Interface} & \textbf{Command}\\
\midrule
Dashboard UI & \file{python dashboard.py} $\rightarrow$ \link{http://localhost:8000}\\
Dashboard alt port & \file{python dashboard.py --port 9000}\\
Chainlit UI & \file{chainlit run ui.py}\\
CLI full run & \file{python step7\_scraper\_with\_fallback.py}\\
Agent self-check & \file{python agent.py}\\
\bottomrule
\end{tabular}
\end{center}

\textbf{Prerequisites for a run:} Ollama must be up at
\link{http://localhost:11434} with the configured model pulled, and
\file{.env} must contain valid GitHub/Gmail credentials.

% ============================================================================
\section{Logs and troubleshooting}
\label{sec:troubleshoot}

\begin{lstlisting}[language=PowerShell, numbers=none]
Get-Content logs\runs-2026-09-18.jsonl   # every run
Get-Content logs\errors.jsonl            # failures + tracebacks
\end{lstlisting}

Common situations:
\begin{itemize}
  \item \textbf{All boards blocked (403/CAPTCHA):} expected from a datacentre
        IP; the router falls back to Gmail automatically and the deck still
        populates from parsed alert emails.
  \item \textbf{Gmail connection hangs:} IMAP now fails fast (15\,s timeout)
        and \tool{asyncio.wait\_for} (60\,s) bounds each tool call.
  \item \textbf{Run is slow:} pure-CPU Ollama inference; a smaller model or an
        NVIDIA-backed Ollama build cuts latency dramatically.
  \item \textbf{\file{watchfiles} fails to install:} harmless (no Rust
        toolchain); uvicorn works without it.
\end{itemize}

% ============================================================================
\section{Current limitations and roadmap}
\label{sec:roadmap}

\subsection{Known limitations}
\begin{itemize}
  \item Some boards (Indeed/LinkedIn/Naukri/Glassdoor/Foundit) bot-wall a plain
        HTTP client; WeWorkRemotely, Remotive, Arbeitnow and Internshala-WFH
        return live listings every run, so the Gmail fallback is the backup for
        a fully-blocked run, not the norm. A Playwright browser-scraper would
        unlock the rest.
  \item Scrapers return page 1 only.
  \item Gmail IMAP is a blunt \tool{FROM sender} search; Google's own alerts
        are loose matches.
  \item Single-turn UI sessions (each message re-runs with system context).
  \item \tool{INDEED\_ACCESS\_TOKEN} (official remote MCP server) is not wired
        up — it needs its own OAuth flow.
\end{itemize}

\subsection{Roadmap}
\begin{itemize}
  \item[\textcolor{green!50!black}{\checkmark}] Graphical dashboard; shared
        \file{agent.py}; resume $\to$ skills $\to$ role suggestions;
        remote/outside-India filtering; multi-board search; flashcard deck
        with Apply buttons; run/error logging with History drawer.
  \item[\textcolor{gray}{$\circ$}] Step 8: scheduled daily digest (Windows Task
        Scheduler) + JSON persistence.
  \item[\textcolor{gray}{$\circ$}] Browser-driven scraping (Playwright).
  \item[\textcolor{gray}{$\circ$}] Resume tailoring: per-job \LaTeX $\to$ PDF.
  \item[\textcolor{gray}{$\circ$}] Multi-turn conversation + persisted chat history.
  \item[\textcolor{gray}{$\circ$}] Scraper pagination and more robust card parsing.
\end{itemize}

\vfill
\begin{center}
{\small\color{gray} End of document.\quad JobScope --- generated \today}
\end{center}

\end{document}